%! Logarithmic Function Manipulation — Unit 5, Lesson 5.4 — Exit Ticket (Answer Key)
\documentclass[10pt]{article}
\usepackage{precalculus-article}
\usepackage{precalculus-key}   % pulls in -boxes; do NOT also load -boxes
\newcommand{\TallMath}[1]{$\displaystyle #1\rule[-1.4em]{0pt}{3.2em}$}

\begin{document}

\pageheader{Unit 5, Lesson 5.4}{Exit Ticket --- Answer Key}
\namedateperiod

\vspace{0.18in}

\begin{enumerate}[leftmargin=1.8em, itemsep=18pt]

\item Rewrite $\log_4(16x^2)$ in the form $a + b\log_4(x)$.

\vspace{6pt}
Show your steps:

\vspace{4pt}
$\log_4(16x^2) = \log_4($ \ans{16} $) + \log_4($ \ans{$x^2$} $)$
\hfill (product property)

\vspace{8pt}
$=$ \ans{2} $+$ \ans{2} $\log_4(x)$
\hfill (evaluate $\log_4(16)$; power property)

\vspace{8pt}
$a = $ \ans{2} \qquad $b = $ \ans{2}

\begin{teachernote}
Common error: students may apply the power property first and write $\log_4(16) + 2\log_4(x)$
before evaluating $\log_4(16) = 2$. Both orders are correct; reward either path if the final
answer is $2 + 2\log_4(x)$. Watch for students who write $b = 4$ (confusing the base with
the exponent) or $a = 16$ (forgetting to evaluate the log).
\end{teachernote}

\vspace{12pt}

\item Use the change-of-base formula to express $\log_7(50)$ as a ratio of common logarithms.

\vspace{6pt}
$\log_7(50) = \dfrac{\log({\color{keyred}\mathbf{50}})}{\log({\color{keyred}\mathbf{7}})}$

\vspace{10pt}
Decimal approximation (optional): $\log_7(50) \approx $ \ans{$2.010$}

\begin{teachernote}
$\log(50)/\log(7) \approx 1.699/0.845 \approx 2.010$. Students who write $\log(50-7)$ or
$\log(50)/7$ are confusing change of base with subtraction or scaling --- address directly.
\end{teachernote}

\end{enumerate}

\end{document}
